\section{Concluding Remarks}
\label{thesis:concl-remarks}

Integrating \acp{TEE} into distributed applications requires careful
consideration of the attacker model and the limitations of this technology.
Simply deploying a workload inside an enclave or confidential virtual machine,
in a ``lift-and-shift'' manner, might create a false sense of security where
users expect that their workload is secure while, in reality, it might still be
exposed to attacks. Looking at the \ac{TEE} threat model, infrastructure
providers are considered untrusted entities and, with this in mind, the term
``confidential computing'' has been adopted in the market to describe workloads
running inside \acp{TEE}, with the promise that providers \emph{cannot} look at
or tamper with applications running on their own infrastructure. 

However, our research shows that, in some cases, cloud providers could actually
break this promise if they wanted to, since they play a significant role in the
provisioning and attestation of \ac{TEE} workloads. Even though there is a good
possibility that this problem will be solved in the future, cloud providers can
still rely on side channels, physical attacks, and potential vulnerabilities in
the \ac{TEE} implementation to achieve their malicious intent. These attacks can
be partially mitigated in software, but researchers have demonstrated that some
of them are still a realistic threat. Thus, it is debatable whether cloud
providers should really be considered untrusted in the \ac{TEE} threat model, or
if they should rather be regarded as ``honest-but-curious'' entities that in
theory have the capabilities to compromise a \ac{TEE}, but in practice will not
perform active attacks. Yet, it is possible that these attacks might also be
carried out by remote adversaries who somehow succeeded in taking control of a
cloud server. Here, users should also trust cloud providers to adopt adequate
security mechanisms to prevent such compromise from happening.

Even considering their limitations, \acp{TEE} significantly increase security by
making it much harder for an adversary to compromise the confidentiality and
integrity of applications. To break into a \ac{TEE} instance, malicious cloud
providers, co-resident tenants or remote attackers would need to perform
sophisticated attacks that would require time, effort and expertise to succeed.
Thus, this technology is an important building block towards making more robust
systems, coming with relatively low cost and overhead. As \acp{TEE} are being
used in production environments today, it is important to provide tools and
support for using this technology correctly, while at the same time raising
awareness on its challenges and limitations.